\documentclass[12pt]{article}
\usepackage[utf8]{inputenc}
\usepackage[T1]{fontenc}
\usepackage{hyperref}
\title{Vafsi Language Resources for K-12 Curriculum Development}
\author{Compiled by Claude Sonnet 5 (claude-sonnet-5) for Andr\'as Kornai}
\date{}
\begin{document}
\maketitle
\section{Introduction}
Vafsi (ISO 639-3: vaf) is a Tatic/Central Plateau Iranian language of Markazi province, Iran, spoken in villages such as Vafs, Chehreqan, Gurchan and Fark. Resources found are academic; no school materials were found.
\section{Primary Educational Resources}
None found.
\section{Linguistic and Reference Materials}
\begin{itemize}
\item Vafsi dialect (Wikipedia): overview, classification and grammar notes, citing Donald Stilo's 1971 dissertation \emph{A Grammar of Vafsi-Tati}; a larger grammar is reported in draft, publication unconfirmed. \url{https://en.wikipedia.org/wiki/Vafsi_dialect}
\item \emph{Vafsi Folk Tales} (Reichert, 2004; ed. Stilo and Marzolph): 24 Gurchani tales with English translation, notes and glossary. Preface: \url{https://reichert-verlag.de/media/pdf/9783895004230_preface.pdf} Sample: \url{https://reichert-verlag.de/media/pdf/9783895004230_sample.pdf}
\item ASJP wordlist for Tatic Vafsi. \url{https://asjp.clld.org/languages/TATIC_VAFSI}
\item ILEX Foundation note on Stilo's Iranian folklore and dialects project. \url{https://ilexfoundation.org/?p=233}
\end{itemize}
No dedicated dictionary was found.
\section{Religious and Community Resources}
None found.
\section{Digital and Audio Resources}
None found beyond the wordlist above.
\section{Curriculum Development Context}
Vafsi is not known to be taught in schools. The folk tales and Stilo's grammar would be the main sources for building native-language reading material, with Persian as the school language.
\end{document}